\section{Introduction}
\label{sec:intro}

The dissertation~\cite{wehrman12} set out to give a local reasoning
principle for programs running on an x86-like machine, where each processor
buffers its writes and other processors cannot see a write until it leaves
the buffer. Its logic has three separating conjunctions. Spatial
conjunction $P \fsep Q$ splits memory by address, as in separation logic.
Temporal conjunction $P \fseq Q$ splits it by time: $P$ describes earlier
writes and $Q$ later ones, possibly to the same addresses. Strong temporal
conjunction $P \fsseq Q$ requires both. Each has a frame rule, and a
concurrency rule and a shared invariant in the style of concurrent
separation logic sit on top.

The dissertation was candid about where it stopped. Its statement of
contributions lists two major components as unfinished and admits a third
gap:
\begin{enumerate}[nosep]
  \item there is no proof system for the assertion language;
  \item the soundness proof for the specification logic is incomplete; and
  \item there are few worked examples, so it is unclear how one would
        actually use the rules to verify a program.
\end{enumerate}
It also hypothesized that the logic was sound and that the remaining proofs
were ``mostly straightforward inductions.''

I have wondered since whether that was true. This report returns to those
three items, fourteen years later, to satisfy that curiosity, with a
mechanization in the Rocq prover. Its findings are as follows.

\paragraph{The soundness theorem does not hold as stated.}
Several auxiliary claims are false, but each can be fixed without
changing the design (Section~\ref{sec:local-errors}). More seriously, the rules themselves are
inconsistent with any adequate semantics: from \rn{skip}, \rn{cons},
\rn{inv} and \rn{frame-tm} alone one can derive that a $\cskip$ changes the
contents of a memory cell (Section~\ref{sec:inconsistent}). No choice of
safety predicate can repair that.

\paragraph{The cause is general.}
An unrestricted temporal frame rule, a right unit for temporal conjunction,
the rule of consequence, and a command that gives away ownership of a
cell cannot coexist once some history that held before the transfer is
false after it (Section~\ref{sec:trilemma}). The observation is
independent of the memory model, and a sequential conjunction in which
later values override earlier ones supplies such histories.

\paragraph{A repaired calculus is sound.}
The repair replaces the invariant rule by lock and unlock rules that
transfer the invariant, restricts the temporal frame rule to commands that do not release the lock, and adjusts four side conditions
(Section~\ref{sec:changes}). The repaired calculus is proved sound by
induction on derivations, with an adequacy theorem for the dissertation's
machine (Section~\ref{sec:results}). The proofs are checked by the Rocq
kernel and use no axioms.

\paragraph{Most of the design survives.}
The machine, the assertion language and its interpretation of atoms, the
unit laws for $\femp$, overlapping temporal history, the temporal and
spatial frame rules, the concurrency rule, and the original load and store
axioms are all retained. So is the dissertation's assertion algebra,
with one family of exceptions: the original laws that rebracket or exchange
through strong temporal conjunction fail, and with them the strong temporal
frame rule. The exchange laws between spatial and temporal conjunction also
fail or are vacuous (Appendix~\ref{app:statements}).

\paragraph{Assertions have a proof system.}
A calculus of 47 rules for entailment between assertions is proved sound
(Section~\ref{sec:assertion-rules}). Its completeness is not established and was not
a goal, but it is enough: the larger derivations below use no semantic
entailment at all.

\paragraph{The logic can be used.}
Section~\ref{sec:proofs} derives a shared container in which a thread
acquires the lock, removes a cell, releases the lock, writes the cell
privately, and later returns it, while another processor's write stays
buffered throughout. It re-derives the dissertation's own message-passing
example. And it builds named locks on top of the single global lock, as
derived rules, and uses one to pass a message of any length from a producer
to a consumer.

\medskip
This report is for a reader who knew the original dissertation but has
forgotten its details. Section~\ref{sec:original} recalls the original model and
calculus in the original notation. Sections~\ref{sec:wrong}
and~\ref{sec:changes} say what was wrong and what changed.
Sections~\ref{sec:model} to~\ref{sec:proofs} present the repaired model,
calculus, results and derivations. Proofs are not reproduced; every claim
cites the name of its Rocq theorem, and the appendix gives fuller
statements.
